%==============================================================================
%  Lab Report Template — Lab 1: Introduction
%  Course : Introduction to Machine Learning (with Python & scikit-learn)
%  Follows: notebooks/01-introduction-code.ipynb
%           "Classifying Iris Species with k-Nearest Neighbors"
%
%  Compile: pdflatex lab1-introduction.tex   (run twice for references)
%       or: latexmk -pdf lab1-introduction.tex
%
%  HOW TO USE
%  ----------
%  1. Fill in the metadata in the "STUDENT INFORMATION" block below.
%  2. Run the notebook, then replace every  % TODO  / \textit{TODO} with the
%     values and plots you obtained.
%  3. Put your figures in a  figures/  subfolder and use \includegraphics.
%  4. Delete the guidance comments (lines starting with %) before submitting.
%==============================================================================
\documentclass[11pt,a4paper]{article}

%------------------------------ packages --------------------------------------
\usepackage[T1]{fontenc}
\usepackage[utf8]{inputenc}
\usepackage{lmodern}
\usepackage[margin=2.5cm]{geometry}
\usepackage{graphicx}
\usepackage{booktabs}
\usepackage{amsmath,amssymb}
\usepackage{xcolor}
\usepackage{listings}
\usepackage{caption}
\usepackage{subcaption}
\usepackage{enumitem}
\usepackage{float}
\usepackage{fancyhdr}
\usepackage[hidelinks]{hyperref}

%------------------------------ metadata --------------------------------------
% ========================= STUDENT INFORMATION ===============================
\newcommand{\Course}{Introduction to Machine Learning}
\newcommand{\LabNo}{Lab 1}
\newcommand{\LabTitle}{Introduction}
\newcommand{\StudentName}{Your Name}          % TODO: your full name
\newcommand{\StudentID}{Student ID}           % TODO: your student number
\newcommand{\DueDate}{2026.09.23}                 % TODO: the due date
\newcommand{\Instructor}{Zhang Li \& Jing Ming}
\newcommand{\Institution}{Shandong Women's University}
% =============================================================================

%------------------------------ colours & code style --------------------------
\definecolor{codebg}{RGB}{245,242,247}
\definecolor{kw}{RGB}{125,15,98}      % keywords (burgundy)
\definecolor{str}{RGB}{31,51,82}      % strings (navy)
\definecolor{cmt}{RGB}{91,100,114}    % comments (muted)
\definecolor{accent}{RGB}{138,28,48}  % headings accent

\lstdefinestyle{py}{
  language=Python,
  backgroundcolor=\color{codebg},
  basicstyle=\ttfamily\small,
  keywordstyle=\color{kw}\bfseries,
  stringstyle=\color{str},
  commentstyle=\color{cmt}\itshape,
  numberstyle=\tiny\color{cmt},
  numbers=left,
  numbersep=8pt,
  frame=leftline,
  framerule=1.2pt,
  rulecolor=\color{accent},
  showstringspaces=false,
  breaklines=true,
  tabsize=4,
  xleftmargin=14pt,
  aboveskip=10pt,
  belowskip=6pt,
}
\lstset{style=py}

%------------------------------ headers/footers -------------------------------
\pagestyle{fancy}
\fancyhf{}
\lhead{\small \Course}
\rhead{\small \LabNo\ — \LabTitle}
\cfoot{\small \StudentName\ (\StudentID)}
\rfoot{\small \thepage}
\renewcommand{\headrulewidth}{0.4pt}
\renewcommand{\footrulewidth}{0.4pt}

%------------------------------ helpers ---------------------------------------
% A grey placeholder box to stand in for a figure you have not added yet.
\newcommand{\placeholder}[2][8cm]{%
  \fbox{\parbox[c][5cm][c]{#1}{\centering\color{cmt}\itshape #2}}%
}

%==============================================================================
\begin{document}

%------------------------------ title block -----------------------------------
\begin{titlepage}
\centering
\vspace*{2cm}
{\LARGE\bfseries \Course\par}
\vspace{0.6cm}
{\color{accent}\rule{\linewidth}{1.5pt}}\par
\vspace{0.5cm}
{\Huge\bfseries \LabNo: \LabTitle\par}
\vspace{0.3cm}
{\large Classifying Iris Species with k-Nearest Neighbors\par}
\vspace{0.4cm}
{\color{accent}\rule{\linewidth}{1.5pt}}\par
\vspace{2.3cm}

\begin{tabular}{rl}
  \textbf{Student:}    & \StudentName \\[4pt]
  \textbf{Student ID:} & \StudentID \\[4pt]
  \textbf{Instructor:} & \Instructor \\[4pt]
  \textbf{Institution:}& \Institution \\[4pt]
  \textbf{Date:}       & \DueDate \\
\end{tabular}

\vfill
{\color{cmt}\small Based on \texttt{notebooks/01-introduction-code.ipynb}.
Replace all placeholder text before submission.}
\end{titlepage}

\tableofcontents
\newpage

%==============================================================================
\section{Objectives}
%------------------------------------------------------------------------------
In this first lab you rebuild the book's very first machine learning
application step by step. By the end you should be able to:
\begin{itemize}[leftmargin=*]
  \item Load a built-in scikit-learn dataset and inspect it.
  \item Split data into \textbf{training} and \textbf{test} sets.
  \item Build a \textbf{k-Nearest Neighbors} classifier.
  \item Use the \lstinline|fit| $\rightarrow$ \lstinline|predict| $\rightarrow$
        \lstinline|score| workflow to train, predict, and evaluate a model.
\end{itemize}

%==============================================================================
\section{Background}
%------------------------------------------------------------------------------
In \emph{supervised learning} a model learns from labelled examples
$(\mathbf{x}_i, y_i)$. Because each iris belongs to one of three species, this
is a \textbf{multiclass classification} problem (the possible answers are
\emph{classes}; the answer for one flower is its \emph{label}).

To estimate how well a model \textbf{generalises} to unseen data, we hold out a
\textbf{test set}: the model is trained only on the training set and evaluated on
the test set. Accuracy summarises performance:
\[
  \text{accuracy} = \frac{\text{\# correct predictions}}{\text{\# samples}}.
\]

\noindent\textbf{k-Nearest Neighbors (k-NN).} To classify a new point, k-NN looks
at the $k$ closest training points and takes a majority vote. With
\lstinline|n_neighbors=1| the prediction is simply the label of the single
nearest training point. Every scikit-learn model follows the same pattern:
\lstinline|fit| (train), \lstinline|predict| (classify new data),
\lstinline|score| (accuracy).

%==============================================================================
\section{Dataset}
%------------------------------------------------------------------------------
We use the classic \textbf{Iris} dataset, bundled with scikit-learn as
\lstinline|load_iris()|. Each flower is described by four measurements, and the
target is one of three species.

\begin{table}[H]
\centering
\caption{Dataset summary (fixed facts, provided for reference).}
\label{tab:dataset}
\begin{tabular}{@{}lr@{}}
\toprule
\textbf{Property} & \textbf{Value} \\
\midrule
Source              & \lstinline|sklearn.datasets.load_iris| \\
Task type           & Multiclass classification (3 classes) \\
Number of samples   & 150 \\
Number of features  & 4 (sepal/petal length \& width, cm) \\
Classes             & setosa, versicolor, virginica (50 each) \\
Split               & \lstinline|train_test_split(..., random_state=0)| \\
Training samples    & 112 \\
Test samples        & 38 \\
\bottomrule
\end{tabular}
\end{table}

% Optional: include the pair plot ("look at your data") if you made one.
\begin{figure}[H]
\centering
% \includegraphics[width=0.75\linewidth]{figures/iris_pairplot.png}
\placeholder[11cm]{(Optional) Pair plot of the four features,\\ coloured by species.}
\caption{Scatter-matrix of the Iris features (optional ``look at your data'' step).}
\label{fig:pairplot}
\end{figure}

%==============================================================================
\section{Methodology}
%------------------------------------------------------------------------------
The notebook is organised as seven short exercises built on this setup:

\begin{lstlisting}[caption={Imports (Setup).},label={lst:setup}]
import numpy as np
from sklearn.datasets import load_iris
from sklearn.model_selection import train_test_split
from sklearn.neighbors import KNeighborsClassifier
\end{lstlisting}

\begin{enumerate}[leftmargin=*]
  \item \textbf{Load the data} into \lstinline|iris| with \lstinline|load_iris()|.
  \item \textbf{Explore} the shape and the per-class counts.
  \item \textbf{Split} into train/test with \lstinline|random_state=0|.
  \item \textbf{Build \& train} \lstinline|KNeighborsClassifier(n_neighbors=1)|.
  \item \textbf{Predict} the species of a new flower \lstinline|[5.0, 2.9, 1.0, 0.2]|.
  \item \textbf{Evaluate} accuracy on the test set with \lstinline|score|.
  \item \textbf{(Bonus)} Vary $k$ from 1 to 10 and compare test accuracy.
\end{enumerate}

%==============================================================================
\section{Experiments and Results}
%------------------------------------------------------------------------------
% Run the notebook and fill in every \textit{TODO} with what you obtained.

\subsection{Exercise 1 — Load the data}
Report what the dataset contains.
\begin{itemize}[leftmargin=*]
  \item Feature names: \textit{TODO}
  \item Target names: \textit{TODO} \quad(expected setosa, versicolor, virginica)
  \item Data shape: \textit{TODO} \quad(expected $(150, 4)$)
\end{itemize}

\subsection{Exercise 2 — Explore the data}
\begin{itemize}[leftmargin=*]
  \item Samples, features: \textit{TODO}
  \item Per-class counts: \textit{TODO} \quad(expected 50 / 50 / 50)
\end{itemize}

\subsection{Exercise 3 — Train/test split}
\begin{itemize}[leftmargin=*]
  \item \lstinline|X_train| shape: \textit{TODO} \quad(expected 112 samples)
  \item \lstinline|X_test| shape: \textit{TODO} \quad(expected 38 samples)
\end{itemize}

\subsection{Exercise 4 — Build and train the classifier}
State the model you trained and confirm it fitted successfully.
\begin{itemize}[leftmargin=*]
  \item Model: \lstinline|KNeighborsClassifier(n_neighbors=1)|
  \item Fitted on: \textit{TODO} training samples
\end{itemize}

\subsection{Exercise 5 — Predict a new flower}
For the new measurements \lstinline|[5.0, 2.9, 1.0, 0.2]|:
\begin{itemize}[leftmargin=*]
  \item Predicted class index: \textit{TODO}
  \item Predicted species name: \textit{TODO} \quad(expected \emph{setosa})
\end{itemize}

\subsection{Exercise 6 — Evaluate on the test set}
\begin{itemize}[leftmargin=*]
  \item Test accuracy: \textit{TODO} \quad(expected $\approx 0.97$, must be $> 0.90$)
\end{itemize}

\subsection{Exercise 7 (bonus) — Does \texorpdfstring{$k$}{k} matter?}
Record the test accuracy for each number of neighbours.
\begin{table}[H]
\centering
\caption{Test accuracy as the number of neighbours $k$ varies.}
\label{tab:kvalues}
\begin{tabular}{@{}cccccccccc@{}}
\toprule
$k$ & 1 & 2 & 3 & 4 & 5 & 6 & 7 & 8 & 9/10 \\
\midrule
Test acc. & \textit{TODO} & \textit{TODO} & \textit{TODO} & \textit{TODO}
          & \textit{TODO} & \textit{TODO} & \textit{TODO} & \textit{TODO} & \textit{TODO} \\
\bottomrule
\end{tabular}
\end{table}
% Optional: plot the table above.
\begin{figure}[H]
\centering
% \includegraphics[width=0.7\linewidth]{figures/k_vs_accuracy.png}
\placeholder[10cm]{(Optional) Plot of test accuracy vs.\ $k$.}
\caption{Test accuracy as a function of the number of neighbours (optional).}
\label{fig:kplot}
\end{figure}

%==============================================================================
\section{Discussion}
%------------------------------------------------------------------------------
Answer with evidence from your results above:
\begin{itemize}[leftmargin=*]
  \item \textbf{Generalisation.} Why do we evaluate on a separate test set instead
        of on the training data?
  \item \textbf{Prediction (Ex.~5).} Does the predicted species for the new flower
        seem reasonable given its small petals?
  \item \textbf{Choice of $k$ (Ex.~7).} Which value(s) of $k$ gave the best test
        accuracy? How does $k$ relate to model complexity (small $k$ vs.\ large $k$)?
  \item \textbf{Limitations.} What are the drawbacks of k-NN noted in the course
        (speed, many features, scaling)?
\end{itemize}
% TODO: write 1--2 paragraphs answering the questions above.

%==============================================================================
\section{Conclusion}
%------------------------------------------------------------------------------
% 3--5 sentences summarising what you learned and your main result.
% TODO.

%==============================================================================
\begin{thebibliography}{9}
\bibitem{mueller2016}
A.~C. M\"uller and S.~Guido,
\emph{Introduction to Machine Learning with Python}.
O'Reilly Media, 2016.

\bibitem{sklearn}
F.~Pedregosa et al.,
``Scikit-learn: Machine Learning in Python,''
\emph{Journal of Machine Learning Research}, vol.~12, pp.~2825--2830, 2011.
\end{thebibliography}

%==============================================================================
\appendix
\section{Source Code}
%------------------------------------------------------------------------------
% Paste your completed code, or reference the notebook you ran.
Full code: \texttt{notebooks/01-introduction-code.ipynb}. Reference solution:

\begin{lstlisting}[caption={End-to-end reference solution.}]
import numpy as np
from sklearn.datasets import load_iris
from sklearn.model_selection import train_test_split
from sklearn.neighbors import KNeighborsClassifier

iris = load_iris()                                        # Ex1
print(iris.data.shape, np.bincount(iris.target))          # Ex2

X_train, X_test, y_train, y_test = train_test_split(      # Ex3
    iris.data, iris.target, random_state=0)

knn = KNeighborsClassifier(n_neighbors=1)                 # Ex4
knn.fit(X_train, y_train)

X_new = np.array([[5.0, 2.9, 1.0, 0.2]])                  # Ex5
print(iris.target_names[knn.predict(X_new)][0])           # -> setosa

print(round(knn.score(X_test, y_test), 3))                # Ex6  (~0.974)

for k in range(1, 11):                                    # Ex7
    m = KNeighborsClassifier(n_neighbors=k).fit(X_train, y_train)
    print(k, round(m.score(X_test, y_test), 3))
\end{lstlisting}

\end{document}
